% GENERATED FILE. Edit canonical lessons, then run npm run generate:books.
% canonical-source-hash: fnv1a64:79edc40f74422dea
% canonical-lessons: AR-C502-tharthar, AR-C502-mutahaffiz, AR-C502-fuduli, AR-C502-hanun, AR-C502-wadud

\chapter{Talkers and the curious}
\label{ch:ar-talkers-and-the-curious}
\hlchaptermodality{502}

\begin{chapteropening}
\textbf{By the end of this chapter:} \emph{I can describe a person as talkative, reserved, nosy, affectionate or friendly.}

Complete the last lesson of chapter 502: I can describe a person as talkative, reserved, nosy, affectionate or friendly.
\end{chapteropening}

\section[\texorpdfstring{\ar{ثرثار} (tharthār)}{ثرثار (tharthār)}]{\ar{ثرثار} (tharthār) — talkative; a chatterbox}

\label{lesson:AR-C502-tharthar}



\begin{quote}
Before the new one: say the Arabic for sensible, then the Arabic for confidence.
\end{quote}

\subsection*{You'll want to know: \ar{ثرثار}}
\textbf{\ar{ثرثار}} — \emph{tharthār} — \textquotedblleft{}talkative; a chatterbox\textquotedblright{}.

\begin{quote}
\ar{جارنا} \ar{ثرثار} \ar{ويتكلم} \ar{ساعات}.
\end{quote}

\begin{quote}
\emph{jārunā tharthār wa-yatakallamu sāʿāt.}
\end{quote}

\begin{quote}
Our neighbour is talkative and talks for hours.
\end{quote}

\begin{cousinweb}
\textbf{\ar{ثرثار}} (\emph{tharthār}) goes with the verb \emph{tharthara}, \textquotedblleft{}to chatter\textquotedblright{}; its doubled shape, \emph{thar-thar}, is usually said to echo endless talk, but the origin is not settled.
\end{cousinweb}

\begin{grammarlens}[title={What you've built}]
One more word for describing what you lived and felt.
\end{grammarlens}

\subsection*{Your turn}
\begin{itemize}
\raggedright
  \item \emph{Say it:} \emph{tharthār}
  \item \emph{Say it:} \emph{tharthār}, once more
  \item \emph{Say it:} \emph{jārunā tharthār wa-yatakallamu sāʿāt.} — Our neighbour is talkative and talks for hours.
  \item \emph{Recall:} say the Arabic for sensible, then the Arabic for confidence, then say \emph{tharthār} again
\end{itemize}

\begin{tcolorbox}[breakable,colback=teal!4,colframe=teal!35!black,title={Before you move on}]
What does \ar{ثرثار} mean? (\textquotedblleft{}Talkative; a chatterbox\textquotedblright{}.) Say it once more.
\end{tcolorbox}
\section[\texorpdfstring{\ar{متحفظ} (mutaḥaffiẓ)}{متحفظ (mutaḥaffiẓ)}]{\ar{متحفظ} (mutaḥaffiẓ) — reserved}

\label{lesson:AR-C502-mutahaffiz}



\begin{quote}
Before the new one: say the Arabic for confidence, then the Arabic for talkative.
\end{quote}

\subsection*{You'll want to know: \ar{متحفظ}}
\textbf{\ar{متحفظ}} — \emph{mutaḥaffiẓ} — \textquotedblleft{}reserved\textquotedblright{}.

\begin{quote}
\ar{زوج} \ar{خالتي} \ar{رجل} \ar{متحفظ} \ar{ولا} \ar{يتكلم} \ar{كثيرا}.
\end{quote}

\begin{quote}
\emph{zawj khālatī rajul mutaḥaffiẓ wa-lā yatakallamu kathīran.}
\end{quote}

\begin{quote}
My aunt's husband is a reserved man and does not talk much.
\end{quote}

\begin{cousinweb}
\textbf{\ar{متحفظ}} (\emph{mutaḥaffiẓ}) is built on the native Arabic root \emph{ḥ-f-ẓ}, \textquotedblleft{}to keep, to guard\textquotedblright{} (Lane). Ḥafiẓa, to keep, which you know, is from the same root; mutaḥaffiẓ is one who keeps himself to himself.
\end{cousinweb}

\begin{grammarlens}[title={What you've built}]
One more word for describing what you lived and felt.
\end{grammarlens}

\subsection*{Your turn}
\begin{itemize}
\raggedright
  \item \emph{Say it:} \emph{mutaḥaffiẓ}
  \item \emph{Say it:} \emph{mutaḥaffiẓ}, once more
  \item \emph{Say it:} \emph{zawj khālatī rajul mutaḥaffiẓ wa-lā yatakallamu kathīran.} — My aunt's husband is a reserved man and does not talk much.
  \item \emph{Recall:} say the Arabic for confidence, then the Arabic for talkative, then say \emph{mutaḥaffiẓ} again
\end{itemize}

\begin{tcolorbox}[breakable,colback=teal!4,colframe=teal!35!black,title={Before you move on}]
What does \ar{متحفظ} mean? (\textquotedblleft{}Reserved\textquotedblright{}.) Say it once more.
\end{tcolorbox}
\section[\texorpdfstring{\ar{فضولي} (fuḍūlī)}{فضولي (fuḍūlī)}]{\ar{فضولي} (fuḍūlī) — nosy, curious}

\label{lesson:AR-C502-fuduli}



\begin{quote}
Before the new one: say the Arabic for talkative, then the Arabic for reserved.
\end{quote}

\subsection*{You'll want to know: \ar{فضولي}}
\textbf{\ar{فضولي}} — \emph{fuḍūlī} — \textquotedblleft{}nosy, curious\textquotedblright{}.

\begin{quote}
\ar{سأل} \ar{الجار} \ar{الفضولي} \ar{عن} \ar{راتبي}.
\end{quote}

\begin{quote}
\emph{saʾala al-jār al-fuḍūlī ʿan rātibī.}
\end{quote}

\begin{quote}
The nosy neighbour asked about my salary.
\end{quote}

It can be mild (curious) or rude (nosy); the noun \ar{فضول} is curiosity.

\begin{cousinweb}
\textbf{\ar{فضولي}} (\emph{fuḍūlī}) is built on the native Arabic root \emph{f-ḍ-l}, \textquotedblleft{}to be over, to be in excess\textquotedblright{} (Lane). Fuḍūl first meant what is more than enough; fuḍūlī is one who busies himself with what is not his.
\end{cousinweb}

\begin{grammarlens}[title={What you've built}]
One more word for describing what you lived and felt.
\end{grammarlens}

\subsection*{Your turn}
\begin{itemize}
\raggedright
  \item \emph{Say it:} \emph{fuḍūlī}
  \item \emph{Say it:} \emph{fuḍūlī}, once more
  \item \emph{Say it:} \emph{saʾala al-jār al-fuḍūlī ʿan rātibī.} — The nosy neighbour asked about my salary.
  \item \emph{Recall:} say the Arabic for talkative, then the Arabic for reserved, then say \emph{fuḍūlī} again
\end{itemize}

\begin{tcolorbox}[breakable,colback=teal!4,colframe=teal!35!black,title={Before you move on}]
What does \ar{فضولي} mean? (\textquotedblleft{}Nosy, curious\textquotedblright{}.) Say it once more.
\end{tcolorbox}
\section[\texorpdfstring{\ar{حنون} (ḥanūn)}{حنون (ḥanūn)}]{\ar{حنون} (ḥanūn) — affectionate, tender}

\label{lesson:AR-C502-hanun}



\begin{quote}
Before the new one: say the Arabic for reserved, then the Arabic for nosy.
\end{quote}

\subsection*{You'll want to know: \ar{حنون}}
\textbf{\ar{حنون}} — \emph{ḥanūn} — \textquotedblleft{}affectionate, tender\textquotedblright{}.

\begin{quote}
\ar{الأم} \ar{الحنون} \ar{تعانق} \ar{طفلها}.
\end{quote}

\begin{quote}
\emph{al-umm al-ḥanūn tuʿāniqu ṭiflahā.}
\end{quote}

\begin{quote}
The tender mother hugs her child.
\end{quote}

It often keeps one form for a man and a woman: \ar{الأم} \ar{الحنون}.

\begin{cousinweb}
\textbf{\ar{حنون}} (\emph{ḥanūn}) is built on the native Arabic root \emph{ḥ-n-n}, \textquotedblleft{}to yearn, to feel tenderness\textquotedblright{} (Lane). The noun ḥanān is tenderness.
\end{cousinweb}

\begin{grammarlens}[title={What you've built}]
One more word for describing what you lived and felt.
\end{grammarlens}

\subsection*{Your turn}
\begin{itemize}
\raggedright
  \item \emph{Say it:} \emph{ḥanūn}
  \item \emph{Say it:} \emph{ḥanūn}, once more
  \item \emph{Say it:} \emph{al-umm al-ḥanūn tuʿāniqu ṭiflahā.} — The tender mother hugs her child.
  \item \emph{Recall:} say the Arabic for reserved, then the Arabic for nosy, then say \emph{ḥanūn} again
\end{itemize}

\begin{tcolorbox}[breakable,colback=teal!4,colframe=teal!35!black,title={Before you move on}]
What does \ar{حنون} mean? (\textquotedblleft{}Affectionate, tender\textquotedblright{}.) Say it once more.
\end{tcolorbox}
\section[\texorpdfstring{\ar{ودود} (wadūd)}{ودود (wadūd)}]{\ar{ودود} (wadūd) — friendly}

\label{lesson:AR-C502-wadud}



\begin{quote}
Before the new one: say the Arabic for nosy, then the Arabic for affectionate.
\end{quote}

\subsection*{You'll want to know: \ar{ودود}}
\textbf{\ar{ودود}} — \emph{wadūd} — \textquotedblleft{}friendly\textquotedblright{}.

\begin{quote}
\ar{الناس} \ar{في} \ar{القرية} \ar{ودودون}.
\end{quote}

\begin{quote}
\emph{an-nās fī al-qarya wadūdūn.}
\end{quote}

\begin{quote}
The people in the village are friendly.
\end{quote}

\begin{cousinweb}
\textbf{\ar{ودود}} (\emph{wadūd}) is built on the native Arabic root \emph{w-d-d}, \textquotedblleft{}to love, to wish\textquotedblright{} (Lane). Wadūd is one full of affection.
\end{cousinweb}

\begin{grammarlens}[title={What you've built}]
One more word for describing what you lived and felt.
\end{grammarlens}

\subsection*{Your turn}
\begin{itemize}
\raggedright
  \item \emph{Say it:} \emph{wadūd}
  \item \emph{Say it:} \emph{wadūd}, once more
  \item \emph{Say it:} \emph{an-nās fī al-qarya wadūdūn.} — The people in the village are friendly.
  \item \emph{Recall:} say the Arabic for nosy, then the Arabic for affectionate, then say \emph{wadūd} again
\end{itemize}

\begin{tcolorbox}[breakable,colback=teal!4,colframe=teal!35!black,title={Before you move on}]
What does \ar{ودود} mean? (\textquotedblleft{}Friendly\textquotedblright{}.) Say it once more.
\end{tcolorbox}
